% Section 5: Discussion
% Implications, limitations, and future work

\section{Discussion}
\label{sec:discussion}

% ============================================================
% 5.1 Implications
% ============================================================
\subsection{Implications}

Our findings have three practical implications:

\textbf{For LLM serving.}
When upgrading model versions, transfer K state but not V.
The addressing geometry is reusable across versions;
the content representations are not.
This could reduce serving latency and memory costs
during model rollout.

\textbf{For distillation.}
Prior work focuses on matching attention patterns (K geometry).
Our results suggest this is misguided:
K geometry is already generic.
Focus instead on V/content pathways,
which carry task-specific knowledge.

\textbf{For model stitching.}
Heterogeneous reassembly (teacher K + student V)
beats both full models.
This demonstrates that addressing geometry is a composable substrate,
opening new directions in modular model construction.

% ============================================================
% 5.2 Limitations
% ============================================================
\subsection{Limitations}

\textbf{Single domain.}
All experiments use one synthetic OOD task.
Second-domain validation (GSM8K/NQ subsets) is needed
to confirm generalization.

\textbf{Small eval set.}
$n = 14$ evaluation samples limits statistical power.
Bootstrap CI95 is wide; larger eval sets ($n \geq 50$)
are needed for definitive conclusions.

\textbf{No true replication.}
The pipeline is deterministic; 3 seeds are identical.
Variance comes from data, not initialization.
This is a methodological limitation, not a bug.

\textbf{Boundary failures.}
1.7B$\to$0.6B shows negative $\dll$ (weak teacher K hurts).
8B$\to$4B shows EM saturation (pure calibration gain).
These boundary cases need further investigation.

% ============================================================
% 5.3 Future Work
% ============================================================
\subsection{Future Work}

\textbf{Second-domain validation.}
Test on GSM8K (math reasoning) and NQ (entity recognition)
to confirm the K/V asymmetry generalizes beyond our synthetic task.

\textbf{Larger eval sets.}
Scale to $n \geq 50$ evaluation samples with per-hop breakdown
(hop-1 through hop-4) to characterize difficulty effects.

\textbf{Theoretical modeling.}
Why does K geometry transfer but V content doesn't?
A theoretical framework could predict transferability
from model architecture and training data.

\textbf{Real-world deployment.}
Measure actual serving latency and memory savings
from K transfer in production systems.

% ============================================================
% 5.4 Conclusion
% ============================================================
\subsection{Conclusion}

We have shown that KV state transfer across LLMs exhibits
a sharp functional asymmetry:
keys (addressing) transfer, while values (content) do not.
This asymmetry holds across 6 model pairs despite
near-perfect linear alignment (CCA $\rho \approx 1$).
The bottleneck is downstream decode,
where content requires weight-parameterized consumption.

Our results establish that KV state transfer is asymmetric by construction:
addressing geometry is composable across models,
while content is weight-bound.
This has implications for serving, distillation,
and modular model construction.
